\documentclass[11pt,a4paper]{article}

\usepackage[utf8]{inputenc}
\usepackage[T1]{fontenc}
\usepackage[spanish]{babel}
\usepackage{amsmath,amssymb,mathtools}
\usepackage{geometry}
\usepackage{lmodern}
\usepackage{microtype}

\geometry{margin=2.4cm}
\setlength{\parindent}{0pt}
\setlength{\parskip}{0.45em}

\newcommand{\R}{\mathbb{R}}
\newcommand{\sen}{\operatorname{sen}}
\newcommand{\grad}{\operatorname{grad}}
\newcommand{\diver}{\operatorname{div}}
\newcommand{\rot}{\operatorname{rot}}

\begin{document}

\begin{center}
{\LARGE\bfseries Análisis Matemático II}\\[0.4em]
{\Large Tema 3: Derivabilidad y diferenciabilidad en funciones de varias variables}\\[0.4em]
{\large Soluciones}
\end{center}

\section{Ejercicio 1}
\subsection{Solución}

Sea:

\[
f(x,y)=3x-x^2y^2+2x^3y.
\]

Derivando respecto de \(x\):

\[
f_x(x,y)
=
3-2xy^2+6x^2y.
\]

Derivando respecto de \(y\):

\[
f_y(x,y)
=
-2x^2y+2x^3.
\]

Por tanto:

\[
\boxed{
f_x(x,y)=3-2xy^2+6x^2y
}
\]

\[
\boxed{
f_y(x,y)=-2x^2y+2x^3
}
\]

\section{Ejercicio 2}
\subsection{Solución}

Sea:

\[
f(x,y)
=
\cos\left(e^{xy^2}\right)
+
\sen\left(\ln(xy)\right).
\]

Para \(xy>0\):

\[
f_x(x,y)
=
-y^2e^{xy^2}
\sen\left(e^{xy^2}\right)
+
\frac{\cos(\ln(xy))}{x}.
\]

Además:

\[
f_y(x,y)
=
-2xye^{xy^2}
\sen\left(e^{xy^2}\right)
+
\frac{\cos(\ln(xy))}{y}.
\]

Por tanto:

\[
\boxed{
f_x(x,y)
=
-y^2e^{xy^2}
\sen\left(e^{xy^2}\right)
+
\frac{\cos(\ln(xy))}{x}
}
\]

\[
\boxed{
f_y(x,y)
=
-2xye^{xy^2}
\sen\left(e^{xy^2}\right)
+
\frac{\cos(\ln(xy))}{y}
}
\]

\section{Ejercicio 3}
\subsection{Solución}

Sea:

\[
f(x,y,z)
=
\frac{3xz}{x+y}.
\]

Derivando:

\[
f_x(x,y,z)
=
\frac{3yz}{(x+y)^2},
\]

\[
f_y(x,y,z)
=
-\frac{3xz}{(x+y)^2},
\]

\[
f_z(x,y,z)
=
\frac{3x}{x+y}.
\]

Por tanto:

\[
\boxed{
f_x=\frac{3yz}{(x+y)^2},
\qquad
f_y=-\frac{3xz}{(x+y)^2},
\qquad
f_z=\frac{3x}{x+y}
}
\]

Estas derivadas solamente están definidas cuando:

\[
x+y\neq0.
\]

Así:

\[
\boxed{
\text{No existen para todo }(x,y,z)\in\R^3.
}
\]

\section{Ejercicio 4}
\subsection{Solución}

Sea:

\[
f(x,y)
=
1-(x-1)^2-(2-y)^2.
\]

Las pendientes en las direcciones \(X\) e \(Y\) son:

\[
f_x(x,y)
=
-2(x-1),
\]

\[
f_y(x,y)
=
-2(y-2).
\]

En \((1,2)\):

\[
f_x(1,2)=0,
\qquad
f_y(1,2)=0.
\]

Por tanto:

\[
\boxed{
f_x(1,2)=0,
\qquad
f_y(1,2)=0
}
\]

\section{Ejercicio 5}
\subsection{Solución}

Sea:

\[
f(x,y)
=
4-x^2-\frac14y^2.
\]

Su gradiente es:

\[
\nabla f(x,y)
=
\left(-2x,-\frac y2\right).
\]

En \((1,2)\):

\[
\nabla f(1,2)=(-2,-1).
\]

La dirección dada es:

\[
u=
\left(
\cos\frac{\pi}{3},
\sen\frac{\pi}{3}
\right)
=
\left(
\frac12,
\frac{\sqrt3}{2}
\right).
\]

Entonces:

\[
D_u f(1,2)
=
\nabla f(1,2)\cdot u
\]

\[
=
(-2,-1)
\cdot
\left(
\frac12,
\frac{\sqrt3}{2}
\right).
\]

Por tanto:

\[
\boxed{
D_u f(1,2)
=
-1-\frac{\sqrt3}{2}
}
\]

\section{Ejercicio 6}
\subsection{Solución}

Sea:

\[
f(x,y,z)
=
xy^2-4x^2y+z^2.
\]

Su gradiente es:

\[
\nabla f(x,y,z)
=
\left(
y^2-8xy,\,
2xy-4x^2,\,
2z
\right).
\]

En el punto \((1,-1,2)\):

\[
\nabla f(1,-1,2)
=
(9,-6,4).
\]

El vector:

\[
v=(6,2,3)
\]

tiene norma:

\[
\|v\|=7.
\]

Así, el vector unitario correspondiente es:

\[
u=
\frac17(6,2,3).
\]

Entonces:

\[
D_u f(1,-1,2)
=
(9,-6,4)
\cdot
\frac17(6,2,3).
\]

Por tanto:

\[
\boxed{
D_u f(1,-1,2)
=
\frac{54}{7}
}
\]

\section{Ejercicio 7}
\subsection{Solución}

Sea:

\[
f(x,y)
=
3xy^2-2y+5x^2y^2.
\]

Las derivadas de primer orden son:

\[
f_x
=
3y^2+10xy^2,
\]

\[
f_y
=
6xy-2+10x^2y.
\]

Por tanto:

\[
f_{xx}
=
10y^2,
\]

\[
f_{xy}
=
6y+20xy,
\]

\[
f_{yx}
=
6y+20xy,
\]

\[
f_{yy}
=
6x+10x^2.
\]

Así:

\[
\boxed{
f_{xx}=10y^2,
\quad
f_{xy}=f_{yx}=6y+20xy,
\quad
f_{yy}=6x+10x^2
}
\]

En \((-1,2)\):

\[
f_{xy}(-1,2)
=
12-40
=
-28,
\]

\[
f_{yx}(-1,2)
=
-28.
\]

Por tanto:

\[
\boxed{
f_{xy}(-1,2)=f_{yx}(-1,2)=-28
}
\]

\section{Ejercicio 8}
\subsection{Solución}

Sea:

\[
z=a^{x^2-y^2},
\qquad
x=e^{2t},
\qquad
y=e^{-t}.
\]

Por la regla de la cadena:

\[
\frac{dz}{dt}
=
\frac{\partial z}{\partial x}
\frac{dx}{dt}
+
\frac{\partial z}{\partial y}
\frac{dy}{dt}.
\]

Se tiene:

\[
\frac{\partial z}{\partial x}
=
2x\ln(a)\,a^{x^2-y^2},
\]

\[
\frac{\partial z}{\partial y}
=
-2y\ln(a)\,a^{x^2-y^2}.
\]

Además:

\[
\frac{dx}{dt}=2e^{2t}=2x,
\qquad
\frac{dy}{dt}=-e^{-t}=-y.
\]

Por tanto:

\[
\boxed{
\frac{dz}{dt}
=
\left(
4e^{4t}
+
2e^{-2t}
\right)
\ln(a)\,
a^{e^{4t}-e^{-2t}}
}
\]

\section{Ejercicio 9}
\subsection{Solución}

Sea:

\[
z=x^2-3xy+2y^2,
\]

\[
x=3\sen(2t),
\qquad
y=4\cos(2t).
\]

Por la regla de la cadena:

\[
\frac{dz}{dt}
=
(2x-3y)\frac{dx}{dt}
+
(-3x+4y)\frac{dy}{dt}.
\]

Como:

\[
\frac{dx}{dt}=6\cos(2t),
\]

\[
\frac{dy}{dt}=-8\sen(2t),
\]

sustituyendo y simplificando:

\[
\boxed{
\frac{dz}{dt}
=
-46\sen(4t)-72\cos(4t)
}
\]

\section{Ejercicio 10}
\subsection{Solución}

Sea:

\[
z=3x^2-2xy+y^2,
\]

\[
x=3u+2v,
\qquad
y=4u-v.
\]

Sustituyendo:

\[
z
=
3(3u+2v)^2
-
2(3u+2v)(4u-v)
+
(4u-v)^2.
\]

Simplificando:

\[
z
=
19u^2+18uv+17v^2.
\]

Por tanto:

\[
\boxed{
\frac{\partial z}{\partial u}
=
38u+18v
}
\]

\[
\boxed{
\frac{\partial z}{\partial v}
=
18u+34v
}
\]

\section{Ejercicio 11}
\subsection{Solución}

Sea:

\[
z=
\frac{2x-y}{x+3y},
\]

\[
x=e^{2u}\cos(3v),
\qquad
y=e^{2u}\sen(3v).
\]

Sustituyendo:

\[
z
=
\frac{
2\cos(3v)-\sen(3v)
}{
\cos(3v)+3\sen(3v)
}.
\]

La variable \(u\) desaparece, luego:

\[
\boxed{
\frac{\partial z}{\partial u}=0
}
\]

Derivando respecto de \(v\):

\[
\boxed{
\frac{\partial z}{\partial v}
=
-\frac{21}
{\left(
\cos(3v)+3\sen(3v)
\right)^2}
}
\]

\section{Ejercicio 12}
\subsection{Solución}

Sea:

\[
f(x,y,z)
=
g(u,v),
\]

con:

\[
u=x^2+2yz,
\qquad
v=y^2+2xz.
\]

Por la regla de la cadena:

\[
f_x
=
2xg_u+2zg_v,
\]

\[
f_y
=
2zg_u+2yg_v,
\]

\[
f_z
=
2yg_u+2xg_v.
\]

Consideramos:

\[
E
=
(y^2-xz)f_x
+
(x^2-yz)f_y
+
(z^2-xy)f_z.
\]

Sustituyendo:

\[
\begin{aligned}
E
={}&
2g_u
\left[
x(y^2-xz)
+
z(x^2-yz)
+
y(z^2-xy)
\right]\\
&+
2g_v
\left[
z(y^2-xz)
+
y(x^2-yz)
+
x(z^2-xy)
\right].
\end{aligned}
\]

Ambos corchetes se anulan:

\[
x y^2-x^2z+x^2z-yz^2+yz^2-xy^2=0,
\]

\[
zy^2-xz^2+x^2y-y^2z+xz^2-x^2y=0.
\]

Por tanto:

\[
\boxed{
(y^2-xz)f_x
+
(x^2-yz)f_y
+
(z^2-xy)f_z
=
0
}
\]

\section{Ejercicio 13}
\subsection{Solución}

Sea:

\[
f(x,y)
=
\begin{cases}
\dfrac{xy^2}{x^2+y^4},
&
(x,y)\neq(0,0),\\[2mm]
0,
&
(x,y)=(0,0).
\end{cases}
\]

Según:

\[
x=0,
\]

se obtiene:

\[
f(0,y)=0.
\]

Según:

\[
x=y^2,
\]

se obtiene:

\[
f(y^2,y)
=
\frac{y^4}{2y^4}
=
\frac12.
\]

Por tanto:

\[
\boxed{
f\text{ no es continua en }(0,0)
}
\]

Las derivadas parciales en el origen son:

\[
f_x(0,0)
=
\lim_{h\to0}
\frac{f(h,0)-f(0,0)}{h}
=
0,
\]

\[
f_y(0,0)
=
\lim_{k\to0}
\frac{f(0,k)-f(0,0)}{k}
=
0.
\]

Así:

\[
\boxed{
f_x(0,0)=f_y(0,0)=0
}
\]

Por tanto, la función es derivable en el sentido de existencia de sus derivadas parciales.

Como toda función diferenciable es continua:

\[
\boxed{
f\text{ no es diferenciable en }(0,0)
}
\]

\section{Ejercicio 14}
\subsection{Solución}

Sea:

\[
f(x,y)
=
\begin{cases}
x\sen\left(
\dfrac1{x^2+y^2}
\right),
&
(x,y)\neq(0,0),\\[2mm]
0,
&
(x,y)=(0,0).
\end{cases}
\]

Se cumple:

\[
|f(x,y)|
\leq
|x|.
\]

Como:

\[
|x|\to0
\]

cuando:

\[
(x,y)\to(0,0),
\]

se obtiene:

\[
\boxed{
f\text{ es continua en }(0,0)
}
\]

Para la derivada parcial respecto de \(x\):

\[
f_x(0,0)
=
\lim_{h\to0}
\frac{h\sen(1/h^2)}{h}
=
\lim_{h\to0}
\sen\left(\frac1{h^2}\right).
\]

Este límite no existe.

Por tanto:

\[
\boxed{
f_x(0,0)\text{ no existe}
}
\]

mientras que:

\[
f_y(0,0)=0.
\]

Así:

\[
\boxed{
f\text{ no es derivable en }(0,0)
}
\]

y, en consecuencia:

\[
\boxed{
f\text{ no es diferenciable en }(0,0)
}
\]

\section{Ejercicio 15}
\subsection{Solución}

Sea:

\[
f(x,y)
=
\begin{cases}
\dfrac{xy}{\sqrt{x^2+y^2}},
&
(x,y)\neq(0,0),\\[2mm]
0,
&
(x,y)=(0,0).
\end{cases}
\]

Utilizando:

\[
2|xy|
\leq
x^2+y^2,
\]

se obtiene:

\[
|f(x,y)|
\leq
\frac12\sqrt{x^2+y^2}.
\]

Por tanto:

\[
\boxed{
f\text{ es continua en }(0,0)
}
\]

Las derivadas parciales son:

\[
f_x(0,0)=0,
\qquad
f_y(0,0)=0.
\]

Así:

\[
\boxed{
f\text{ es derivable en }(0,0)
}
\]

Para estudiar la diferenciabilidad:

\[
\frac{
f(x,y)-f(0,0)-f_x(0,0)x-f_y(0,0)y
}{
\sqrt{x^2+y^2}
}
=
\frac{xy}{x^2+y^2}.
\]

Según:

\[
y=x,
\]

se obtiene:

\[
\frac{x^2}{2x^2}
=
\frac12.
\]

Por tanto:

\[
\boxed{
f\text{ no es diferenciable en }(0,0)
}
\]

\section{Ejercicio 16}
\subsection{Solución}

Sea:

\[
f(x,y)
=
\begin{cases}
(x^2+y^2)
\sen\left(
\dfrac1{\sqrt{x^2+y^2}}
\right),
&
(x,y)\neq(0,0),\\[2mm]
0,
&
(x,y)=(0,0).
\end{cases}
\]

Se tiene:

\[
|f(x,y)|
\leq
x^2+y^2.
\]

Por tanto:

\[
\boxed{
f\text{ es continua en }(0,0)
}
\]

Además:

\[
f_x(0,0)=0,
\qquad
f_y(0,0)=0.
\]

Para la diferenciabilidad:

\[
\left|
\frac{f(x,y)}
{\sqrt{x^2+y^2}}
\right|
\leq
\sqrt{x^2+y^2}.
\]

El miembro derecho tiende a \(0\), luego:

\[
\boxed{
f\text{ es diferenciable en }(0,0)
}
\]

\section{Ejercicio 17}
\subsection{Solución}

Sea:

\[
f(x,y)=e^{3x+2y}.
\]

Es composición de funciones continuas y derivables, por lo que es continua y diferenciable en todo \(\R^2\).

Sus derivadas parciales son:

\[
f_x(x,y)
=
3e^{3x+2y},
\]

\[
f_y(x,y)
=
2e^{3x+2y}.
\]

En el origen:

\[
f(0,0)=1,
\]

\[
f_x(0,0)=3,
\qquad
f_y(0,0)=2.
\]

Por tanto:

\[
\boxed{
f\text{ es continua, derivable y diferenciable en }(0,0)
}
\]

\section{Ejercicio 18}
\subsection{Solución}

Sea:

\[
f(x,y)
=
\begin{cases}
\dfrac{xy(x^2-y^2)}{x^2+y^2},
&
(x,y)\neq(0,0),\\[2mm]
0,
&
(x,y)=(0,0).
\end{cases}
\]

Se tiene:

\[
|x^2-y^2|
\leq
x^2+y^2,
\]

por lo que:

\[
|f(x,y)|
\leq
|xy|
\leq
\frac{x^2+y^2}{2}.
\]

Así:

\[
\boxed{
f\text{ es continua en }(0,0)
}
\]

Además:

\[
f_x(0,0)=0,
\qquad
f_y(0,0)=0.
\]

Para la diferenciabilidad:

\[
\left|
\frac{f(x,y)}
{\sqrt{x^2+y^2}}
\right|
\leq
\frac12\sqrt{x^2+y^2}.
\]

Por tanto:

\[
\boxed{
f\text{ es diferenciable en }(0,0)
}
\]

Para \((x,y)\neq(0,0)\):

\[
f_{xy}(x,y)
=
f_{yx}(x,y)
=
\frac{
(x^2-y^2)
\left(
x^4+10x^2y^2+y^4
\right)
}{
(x^2+y^2)^3
}.
\]

En el origen calculamos por definición.

Como:

\[
f_x(0,y)=-y,
\]

se tiene:

\[
f_{xy}(0,0)
=
\lim_{y\to0}
\frac{-y}{y}
=
-1.
\]

Por otra parte:

\[
f_y(x,0)=x,
\]

por lo que:

\[
f_{yx}(0,0)
=
\lim_{x\to0}
\frac{x}{x}
=
1.
\]

Así:

\[
\boxed{
f_{xy}(0,0)=-1,
\qquad
f_{yx}(0,0)=1
}
\]

\section{Ejercicio 19}
\subsection{Solución}

Sea:

\[
f(x,y)
=
\begin{cases}
\cos x\,\sen y+\dfrac{x^2}{y},
&
y\neq0,\\[2mm]
0,
&
y=0.
\end{cases}
\]

En el origen:

\[
f(0,0)=0.
\]

Además:

\[
f_x(0,0)
=
\lim_{h\to0}
\frac{f(h,0)}{h}
=
0,
\]

\[
f_y(0,0)
=
\lim_{k\to0}
\frac{f(0,k)}{k}
=
\lim_{k\to0}
\frac{\sen k}{k}
=
1.
\]

Por tanto, la aproximación lineal sería:

\[
\boxed{
L(x,y)=y
}
\]

Sin embargo, según:

\[
y=x^2,
\]

se tiene:

\[
f(x,x^2)
=
\cos x\,\sen(x^2)+1
\longrightarrow
1.
\]

Por tanto, la función no es siquiera continua en el origen.

Así:

\[
\boxed{
L(x,y)=y\text{ no es una buena aproximación lineal de }f
\text{ cerca de }(0,0)
}
\]

\section{Ejercicio 20}
\subsection{Solución}

Sea:

\[
f(x,y)
=
y\ln x+xy^2.
\]

Su gradiente es:

\[
\nabla f(x,y)
=
\left(
\frac yx+y^2,\,
\ln x+2xy
\right).
\]

En \((1,2)\):

\[
\nabla f(1,2)
=
(6,4).
\]

Por tanto:

\[
\boxed{
\nabla f(x,y)
=
\left(
\frac yx+y^2,\,
\ln x+2xy
\right)
}
\]

\[
\boxed{
\nabla f(1,2)=(6,4)
}
\]

\section{Ejercicio 21}
\subsection{Solución}

La temperatura es:

\[
f(x,y)
=
20-4x^2-y^2.
\]

Su gradiente es:

\[
\nabla f(x,y)
=
(-8x,-2y).
\]

En \((2,-3)\):

\[
\nabla f(2,-3)
=
(-16,6).
\]

La dirección de máximo incremento es la del gradiente.

Normalizando:

\[
\frac{(-16,6)}
{\sqrt{(-16)^2+6^2}}
=
\frac{(-8,3)}{\sqrt{73}}.
\]

Por tanto:

\[
\boxed{
u_{\max}
=
\frac1{\sqrt{73}}(-8,3)
}
\]

\section{Ejercicio 22}
\subsection{Solución}

Sea:

\[
f(x,y)
=
xy\ln(x-1)+y^2.
\]

Su gradiente es:

\[
\nabla f(x,y)
=
\left(
y\ln(x-1)+\frac{xy}{x-1},
\,
x\ln(x-1)+2y
\right).
\]

En \((2,1)\):

\[
\nabla f(2,1)
=
(2,2).
\]

La dirección de máximo incremento es:

\[
\boxed{
u_{\max}
=
\frac1{\sqrt2}(1,1)
}
\]

La dirección de mínimo incremento es:

\[
\boxed{
u_{\min}
=
-\frac1{\sqrt2}(1,1)
}
\]

Las direcciones de incremento nulo son las unitarias ortogonales al gradiente:

\[
\boxed{
u_0
=
\pm
\frac1{\sqrt2}(1,-1)
}
\]

\section{Ejercicio 23}
\subsection{Solución}

La superficie es:

\[
z=x^2+y^3.
\]

En \((3,1)\):

\[
f_x(3,1)=6,
\qquad
f_y(3,1)=3.
\]

El plano tangente es:

\[
z-10
=
6(x-3)
+
3(y-1).
\]

Por tanto:

\[
\boxed{
z
=
6x+3y-11
}
\]

\section{Ejercicio 24}
\subsection{Solución}

Sea:

\[
F(x,y,z)
=
(x-1)^2+(y-2)^2+2z^2-1.
\]

La superficie viene dada por:

\[
F(x,y,z)=0.
\]

Su gradiente es:

\[
\nabla F
=
\left(
2(x-1),
2(y-2),
4z
\right).
\]

En:

\[
P=
\left(
\frac32,
2,
\frac{\sqrt6}{4}
\right),
\]

se obtiene:

\[
\nabla F(P)
=
(1,0,\sqrt6).
\]

Por tanto, el plano tangente es:

\[
(x-\tfrac32)
+
\sqrt6
\left(
z-\frac{\sqrt6}{4}
\right)
=
0.
\]

Simplificando:

\[
\boxed{
x+\sqrt6\,z-3=0
}
\]

\section{Ejercicio 25}
\subsection{Solución}

Sea:

\[
F(x,y,z)
=
3x^2e^{2(y+1)}
+
4\sen(x+2z)
+
z^2
-
13.
\]

La superficie es:

\[
F(x,y,z)=0.
\]

Calculamos:

\[
F_x
=
6xe^{2(y+1)}
+
4\cos(x+2z),
\]

\[
F_y
=
6x^2e^{2(y+1)},
\]

\[
F_z
=
8\cos(x+2z)+2z.
\]

En:

\[
P=(2,-1,-1),
\]

se obtiene:

\[
\nabla F(P)
=
(16,24,6).
\]

El plano tangente es:

\[
16(x-2)
+
24(y+1)
+
6(z+1)
=
0.
\]

Por tanto:

\[
\boxed{
8x+12y+3z-1=0
}
\]

La recta normal tiene como vector director:

\[
(16,24,6),
\]

que podemos simplificar a:

\[
(8,12,3).
\]

Así:

\[
\boxed{
\begin{cases}
x=2+8t,\\
y=-1+12t,\\
z=-1+3t.
\end{cases}
}
\]

\section{Ejercicio 26}
\subsection{Solución}

La superficie indicada en el enunciado es:

\[
F(x,y,z)
=
x^3y^3
+
xyz\cos z
-
xz
+
y
+
x
=
0.
\]

Sin embargo, en el punto indicado:

\[
P=(1,-1,0),
\]

se obtiene:

\[
F(1,-1,0)
=
-1.
\]

Por tanto:

\[
\boxed{
(1,-1,0)\text{ no pertenece a la superficie dada}
}
\]

y, con el enunciado tal como aparece en el documento:

\[
\boxed{
\text{no están definidos el plano tangente ni la recta normal a esa superficie en }(1,-1,0).
}
\]

\section{Ejercicio 27}
\subsection{Solución}

Sea:

\[
F(x,y,z)
=
\left(
x^2y^3-z^4,\,
4x^5y^2z,\,
-y^4z^6
\right).
\]

La divergencia es:

\[
\diver F
=
\frac{\partial F_1}{\partial x}
+
\frac{\partial F_2}{\partial y}
+
\frac{\partial F_3}{\partial z}.
\]

Por tanto:

\[
\boxed{
\diver F
=
2xy^3
+
8x^5yz
-
6y^4z^5
}
\]

El rotacional es:

\[
\rot F
=
\left(
\frac{\partial F_3}{\partial y}
-
\frac{\partial F_2}{\partial z},
\,
\frac{\partial F_1}{\partial z}
-
\frac{\partial F_3}{\partial x},
\,
\frac{\partial F_2}{\partial x}
-
\frac{\partial F_1}{\partial y}
\right).
\]

Así:

\[
\boxed{
\rot F
=
\left(
-4y^3z^6-4x^5y^2,\,
-4z^3,\,
20x^4y^2z-3x^2y^2
\right)
}
\]

\section{Ejercicio 28}
\subsection{Solución}

Sea:

\[
\nabla f
=
(f_x,f_y,f_z).
\]

Entonces:

\[
\rot(\grad f)
=
\left(
f_{zy}-f_{yz},
f_{xz}-f_{zx},
f_{yx}-f_{xy}
\right).
\]

Como las derivadas segundas son continuas, por el teorema de Schwarz:

\[
f_{zy}=f_{yz},
\qquad
f_{xz}=f_{zx},
\qquad
f_{yx}=f_{xy}.
\]

Por tanto:

\[
\boxed{
\rot(\grad f)=0
}
\]

Sea ahora:

\[
F=(F_1,F_2,F_3).
\]

Se tiene:

\[
\rot F
=
\left(
F_{3y}-F_{2z},
F_{1z}-F_{3x},
F_{2x}-F_{1y}
\right).
\]

Entonces:

\[
\begin{aligned}
\diver(\rot F)
={}&
F_{3yx}-F_{2zx}
+
F_{1zy}-F_{3xy}\\
&+
F_{2xz}-F_{1yz}.
\end{aligned}
\]

Aplicando Schwarz:

\[
F_{3yx}=F_{3xy},
\]

\[
F_{2zx}=F_{2xz},
\]

\[
F_{1zy}=F_{1yz}.
\]

Por tanto:

\[
\boxed{
\diver(\rot F)=0
}
\]

\section{Ejercicio 29}
\subsection{Solución}

Para:

\[
f(x,y,z)
=
(xy,y+z^2),
\]

la matriz jacobiana es:

\[
\boxed{
J_f(x,y,z)
=
\begin{pmatrix}
y&x&0\\
0&1&2z
\end{pmatrix}
}
\]

Para:

\[
g(x,y,z)
=
\left(
x^2\cos y,
e^{x+z^2}
\right),
\]

se obtiene:

\[
\boxed{
J_g(x,y,z)
=
\begin{pmatrix}
2x\cos y&-x^2\sen y&0\\
e^{x+z^2}&0&2ze^{x+z^2}
\end{pmatrix}
}
\]

\section{Ejercicio 30}
\subsection{Solución}

Sea:

\[
f(x,y)
=
\left(
x^3y^2-5x^2y^2,\,
y^6-3xy^3+7
\right).
\]

Su matriz jacobiana es:

\[
J_f(x,y)
=
\begin{pmatrix}
3x^2y^2-10xy^2&
2x^3y-10x^2y\\
-3y^3&
6y^5-9xy^2
\end{pmatrix}.
\]

En \((2,-1)\):

\[
\boxed{
J_f(2,-1)
=
\begin{pmatrix}
-8&24\\
3&-24
\end{pmatrix}
}
\]

\section{Ejercicio 31}
\subsection{Solución}

Sea:

\[
f(x,y,z)
=
\left(
z\tan(x^2-y^2),
xy\ln\left(\frac z2\right)
\right).
\]

La matriz jacobiana es:

\[
J_f(x,y,z)
=
\begin{pmatrix}
2xz\sec^2(x^2-y^2)&
-2yz\sec^2(x^2-y^2)&
\tan(x^2-y^2)\\[1mm]
y\ln(z/2)&
x\ln(z/2)&
\dfrac{xy}{z}
\end{pmatrix}.
\]

En:

\[
(2,-2,2),
\]

se tiene:

\[
x^2-y^2=0,
\qquad
\ln(z/2)=0.
\]

Por tanto:

\[
\boxed{
J_f(2,-2,2)
=
\begin{pmatrix}
8&8&0\\
0&0&-2
\end{pmatrix}
}
\]

\section{Ejercicio 32}
\subsection{Solución}

Sea:

\[
f(x,y)
=
(x^2-y^2,2xy).
\]

Su matriz jacobiana es:

\[
J_f(x,y)
=
\begin{pmatrix}
2x&-2y\\
2y&2x
\end{pmatrix}.
\]

El jacobiano es:

\[
|J_f(x,y)|
=
4x^2+4y^2.
\]

Por tanto:

\[
\boxed{
|J_f(x,y)|
=
4(x^2+y^2)
}
\]

\section{Ejercicio 33}
\subsection{Solución}

Para:

\[
f(x,y)=x^2-xy+y^3,
\]

se tiene:

\[
f_{xx}=2,
\qquad
f_{xy}=f_{yx}=-1,
\qquad
f_{yy}=6y.
\]

Por tanto:

\[
\boxed{
H_f(x,y)
=
\begin{pmatrix}
2&-1\\
-1&6y
\end{pmatrix}
}
\]

Para:

\[
g(x,y,z)
=
x^2-xy^2+z^2,
\]

las derivadas segundas son:

\[
g_{xx}=2,
\qquad
g_{xy}=g_{yx}=-2y,
\qquad
g_{xz}=g_{zx}=0,
\]

\[
g_{yy}=-2x,
\qquad
g_{yz}=g_{zy}=0,
\qquad
g_{zz}=2.
\]

Por tanto:

\[
\boxed{
H_g(x,y,z)
=
\begin{pmatrix}
2&-2y&0\\
-2y&-2x&0\\
0&0&2
\end{pmatrix}
}
\]

\section{Ejercicio 34}
\subsection{Solución}

Sea:

\[
f(x,y)
=
y^4+x^3+3x^2+4y^2-4xy-5y+8.
\]

Se tiene:

\[
f_{xx}=6x+6,
\]

\[
f_{xy}=f_{yx}=-4,
\]

\[
f_{yy}=12y^2+8.
\]

En \((1,0)\):

\[
\boxed{
H_f(1,0)
=
\begin{pmatrix}
12&-4\\
-4&8
\end{pmatrix}
}
\]

\section{Ejercicio 35}
\subsection{Solución}

Sea:

\[
f(x,y)
=
e^{y\ln x}
=
x^y.
\]

Sus derivadas segundas son:

\[
f_{xx}
=
y(y-1)x^{y-2},
\]

\[
f_{xy}
=
x^{y-1}
+
yx^{y-1}\ln x,
\]

\[
f_{yx}=f_{xy},
\]

\[
f_{yy}
=
x^y(\ln x)^2.
\]

En \((1,1)\):

\[
f_{xx}(1,1)=0,
\]

\[
f_{xy}(1,1)=f_{yx}(1,1)=1,
\]

\[
f_{yy}(1,1)=0.
\]

Por tanto:

\[
\boxed{
H_f(1,1)
=
\begin{pmatrix}
0&1\\
1&0
\end{pmatrix}
}
\]

\section{Ejercicio 36}
\subsection{Solución}

Sea:

\[
f(x,y)
=
x^2-y^2+2xy.
\]

Su matriz hessiana es:

\[
H_f(x,y)
=
\begin{pmatrix}
2&2\\
2&-2
\end{pmatrix}.
\]

El hessiano es:

\[
|H_f|
=
2(-2)-2\cdot2
=
-8.
\]

Por tanto:

\[
\boxed{
|H_f(2,1)|=-8
}
\]

\section{Ejercicio 37}
\subsection{Solución}

Sean:

\[
f(x,y)
=
\left(
xy,
\sqrt{x^2+y^2}
\right)
\]

y:

\[
g(u,v)
=
\begin{cases}
\dfrac{u^2}{v},&v\neq0,\\[2mm]
0,&v=0.
\end{cases}
\]

Para \(f\):

\[
J_f(x,y)
=
\begin{pmatrix}
y&x\\[1mm]
\dfrac{x}{\sqrt{x^2+y^2}}&
\dfrac{y}{\sqrt{x^2+y^2}}
\end{pmatrix}.
\]

En \((1,1)\):

\[
\boxed{
J_f(1,1)
=
\begin{pmatrix}
1&1\\
\dfrac1{\sqrt2}&\dfrac1{\sqrt2}
\end{pmatrix}
}
\]

Como:

\[
f(1,1)
=
(1,\sqrt2),
\]

y \(\sqrt2\neq0\), utilizamos:

\[
g(u,v)=\frac{u^2}{v}.
\]

Entonces:

\[
J_g(u,v)
=
\begin{pmatrix}
\dfrac{2u}{v}&-\dfrac{u^2}{v^2}
\end{pmatrix}.
\]

Así:

\[
\boxed{
J_g(1,\sqrt2)
=
\begin{pmatrix}
\sqrt2&-\dfrac12
\end{pmatrix}
}
\]

La composición es:

\[
h(x,y)
=
g(f(x,y))
=
\frac{x^2y^2}{\sqrt{x^2+y^2}}.
\]

En \((1,1)\):

\[
\boxed{
J_h(1,1)
=
\begin{pmatrix}
\dfrac{3}{2\sqrt2}&
\dfrac{3}{2\sqrt2}
\end{pmatrix}
}
\]

Además:

\[
J_g(f(1,1))J_f(1,1)
=
\begin{pmatrix}
\sqrt2&-\dfrac12
\end{pmatrix}
\begin{pmatrix}
1&1\\
\dfrac1{\sqrt2}&\dfrac1{\sqrt2}
\end{pmatrix}
\]

\[
=
\begin{pmatrix}
\dfrac{3}{2\sqrt2}&
\dfrac{3}{2\sqrt2}
\end{pmatrix}.
\]

Por tanto:

\[
\boxed{
J_h(1,1)
=
J_g(f(1,1))J_f(1,1)
}
\]

\section{Ejercicio 38}
\subsection{Solución}

Sea:

\[
f(x,y)
=
(f_1(x,y),f_2(x,y))
=
\left(
xy,
\sqrt{x^2+y^2}
\right).
\]

Para:

\[
f_1(x,y)=xy,
\]

la diferencial total es:

\[
\boxed{
df_1
=
y\,dx+x\,dy
}
\]

Para:

\[
f_2(x,y)
=
\sqrt{x^2+y^2},
\]

se tiene:

\[
\boxed{
df_2
=
\frac{x}{\sqrt{x^2+y^2}}\,dx
+
\frac{y}{\sqrt{x^2+y^2}}\,dy
}
\]

En \((1,1)\):

\[
\boxed{
df_1(1,1)
=
dx+dy
}
\]

\[
\boxed{
df_2(1,1)
=
\frac1{\sqrt2}dx
+
\frac1{\sqrt2}dy
}
\]

La diferencial de la función vectorial coincide con su matriz jacobiana:

\[
\boxed{
Df(1,1)
=
\begin{pmatrix}
1&1\\
\dfrac1{\sqrt2}&\dfrac1{\sqrt2}
\end{pmatrix}
}
\]

\end{document}
